\chapter{German Measles (Rubella)}
\index{German Measles (Rubella)}

\section*{Definition and Overview}
German measles, also known as rubella, is a contagious viral infection that causes a mild illness with a distinctive rash, fever, and swollen lymph glands. It is caused by the rubella virus and is spread through droplets from coughing, sneezing, or talking. For most children and adults, rubella is a brief and mild illness. The infection is dangerous, however, if a woman catches it during early pregnancy, because the virus can cross the placenta and cause serious birth defects known as congenital rubella syndrome, including deafness, cataracts, heart defects, and brain damage. Rubella is now rare in many countries because of high vaccination coverage with the measles-mumps-rubella (MMR) vaccine, which provides lifelong protection.

\section*{Epidemiology}
Rubella was once a common childhood infection worldwide, but vaccination has greatly reduced transmission. In countries with high vaccination coverage, cases are now uncommon, although outbreaks can occur when coverage falls. Rubella still circulates in areas with limited vaccination, and infection during pregnancy can cause congenital rubella syndrome. Maintaining high vaccination coverage and protecting pregnant people from exposure remain essential.

\section*{Aetiology and Pathophysiology}
Rubella is caused by infection with the rubella virus, a member of the togavirus family.

\begin{itemize}
    \item \textbf{Transmission:} via respiratory droplets and direct contact with an infected person; the virus enters through the nose and throat
    \item \textbf{Incubation:} symptoms appear 14--21 days after exposure, usually around 16--18 days
    \item \textbf{Infectious period:} from about one week before to at least four days after the rash appears
    \item \textbf{Immune response:} infection produces lifelong immunity; the MMR vaccine provides durable protection
    \item \textbf{Congenital infection:} the virus crosses the placenta in early pregnancy, infecting the fetus and interfering with organ development
    \item \textbf{Pathology:} in the fetus, rubella damages developing organs, particularly the eyes, ears, heart, and brain
\end{itemize}

\section*{Clinical Features}
Rubella is often so mild that it goes unnoticed, especially in children. When symptoms occur, they include:

\begin{itemize}
    \item A pink or red rash that starts on the face and spreads downward to the trunk and limbs, lasting about three days
    \item Low-grade fever
    \item Swollen, tender lymph nodes, particularly behind the ears and at the back of the neck
    \item A runny or stuffy nose, sore throat, and cough in some cases
    \item Joint pain and swelling, especially in adult women, which can persist for weeks
    \item Headache and a general feeling of being unwell
    \item Up to half of infections cause no symptoms at all
\end{itemize}

\section*{Differential Diagnosis}
The rash of rubella resembles several other childhood infections, and laboratory confirmation is often needed.

\begin{itemize}
    \item Measles, which causes a more severe illness with high fever, cough, coryza, and conjunctivitis, and a blotchy rash
    \item Scarlet fever (group A streptococcal infection), which causes a fine rash with a sore throat
    \item Roseola infantum, which affects young infants and causes fever followed by rash
    \item Parvovirus B19 (fifth disease), which causes a characteristic "slapped cheek" rash
    \item Dengue fever and other viral infections with rash
    \item Drug reactions causing a measles-like rash
\end{itemize}

\section*{When to Seek Medical Care}
People with a rubella-like illness should see a doctor for diagnosis, especially women of childbearing age, who need urgent testing if they may be pregnant or exposed. Pregnant women who have been in contact with someone with rubella, or who develop a rash and fever, should contact their doctor or maternity team the same day.

\textbf{Contact your local emergency services for:}
\begin{itemize}
    \item This infection is rarely an emergency in otherwise healthy people; however, any person with a rash who is very unwell, drowsy, or struggling to breathe should be assessed urgently
    \item A pregnant woman with known or suspected rubella exposure needs urgent (not emergency) medical advice from her maternity team
\end{itemize}

\section*{Diagnosis and Investigations}
Because the illness is mild but the consequences in pregnancy are serious, laboratory confirmation is important.

\begin{itemize}
    \item \textbf{Clinical suspicion:} the combination of rash, fever, and lymphadenopathy in a person who may be unvaccinated
    \item \textbf{Blood tests:} rubella IgM antibodies indicate recent infection; rubella IgG antibodies indicate past infection or vaccination
    \item \textbf{PCR testing:} a throat or nasal swab can detect the virus in the first few days of illness
    \item \textbf{Prenatal testing:} if a pregnant woman is infected, amniocentesis may be offered to check whether the fetus is affected
    \item \textbf{Notification:} rubella is a notifiable disease in many countries, and cases are reported to public health authorities
\end{itemize}

\section*{Management}
There is no specific antiviral treatment for rubella; care is supportive. The priority is preventing infection in pregnant women.

\begin{itemize}
    \item \textbf{Symptom relief:} paracetamol for fever and aches, rest, and adequate fluids
    \item \textbf{Isolation:} staying away from work, school, and childcare until at least four days after the rash appears
    \item \textbf{Avoiding pregnant contacts:} keeping infected people away from pregnant women is critical
    \item \textbf{Management in pregnancy:} specialist obstetric and infectious disease input; the risk of fetal damage is highest in the first 12 weeks
    \item \textbf{Immunoglobulin:} in some exposed pregnant women, antibodies may be given, though they do not always prevent infection
    \item \textbf{Public health action:} tracing contacts and identifying susceptible pregnant women for follow-up
\end{itemize}

\section*{Complications}
\begin{itemize}
    \item Congenital rubella syndrome: deafness, cataracts, heart defects, intellectual disability, and growth problems in babies infected in early pregnancy
    \item Miscarriage and stillbirth following infection in early pregnancy
    \item Arthritis and arthralgia, especially in adult women
    \item Rarely, thrombocytopenia (low platelets) causing bruising and bleeding
    \item Rarely, encephalitis, an inflammation of the brain
    \item The risk of severe fetal damage is very high (up to 85\%) when infection occurs in the first 12 weeks of pregnancy
\end{itemize}

\section*{Prognosis and Outcomes}
For children and healthy adults, rubella is a mild, self-limiting illness that resolves within a week and causes no lasting harm. The serious burden of disease falls on unborn babies: infection in early pregnancy frequently causes permanent disability or fetal loss. Because vaccination has eliminated rubella in many countries, congenital rubella syndrome is now very rare here. The prognosis for the community depends on maintaining high vaccination coverage so the virus cannot circulate.

\section*{Prevention and Self-Care}
\begin{itemize}
    \item Vaccination with the MMR vaccine is the cornerstone of prevention: two doses provide lifelong protection
    \item Ensure children receive the MMR vaccine at 12 months and 18 months under the National Immunisation Program
    \item Adults who are unsure of their vaccination history should check with their GP, and women planning pregnancy should have their immunity confirmed
    \item Women should avoid becoming pregnant for at least 28 days after receiving the MMR vaccine
    \item Pregnant women should avoid contact with anyone who may have rubella
    \item People with rubella should stay home from school, work, and childcare until at least four days after the rash begins
    \item Seek antenatal blood tests, which include rubella immunity screening
\end{itemize}

\section*{Sources and Further Reading}
The following reputable sources provide further information for healthcare students and patients seeking reliable, current advice about German measles (rubella).

\begin{itemize}
    \item Australian Government Department of Health and Aged Care. \textit{Rubella (German measles) fact sheet.} https://www.health.gov.au/
    \item Healthdirect Australia. \textit{Rubella (German measles).} https://www.healthdirect.gov.au/
    \item National Centre for Immunisation Research and Surveillance. \textit{Rubella and the MMR vaccine.} https://www.ncirs.org.au/
    \item NHS. \textit{Rubella (German measles).} https://www.nhs.uk/
    \item Centers for Disease Control and Prevention. \textit{Rubella fact sheet.} https://www.cdc.gov/
    \item World Health Organization. \textit{Rubella fact sheet.} https://www.who.int/
\end{itemize}
